% --- LaTeX Homework Template - S. Venkatraman ---

% --- Set document class and font size ---

\documentclass[letterpaper, 11pt]{article}

% --- Package imports ---

\usepackage{
  amsmath, amsthm, amssymb, mathtools, dsfont,	  % Math typesetting
  graphicx, wrapfig, subfig, float,                  % Figures and graphics formatting
  listings, color, inconsolata, pythonhighlight,     % Code formatting
  fancyhdr, sectsty, hyperref, enumerate, enumitem } % Headers/footers, section fonts, links, lists

% --- Page layout settings ---

% Set page margins
\usepackage[left=1.35in, right=1.35in, bottom=1in, top=1.1in, headsep=0.2in]{geometry}

% Anchor footnotes to the bottom of the page
\usepackage[bottom]{footmisc}

% Set line spacing
\renewcommand{\baselinestretch}{1.2}

% Set spacing between paragraphs
\setlength{\parskip}{1.5mm}

% Allow multi-line equations to break onto the next page
\allowdisplaybreaks

% Enumerated lists: make numbers flush left, with parentheses around them
\setlist[enumerate]{wide=0pt, leftmargin=21pt, labelwidth=0pt, align=left}
\setenumerate[1]{label={(\arabic*)}}

% --- Page formatting settings ---

% Set link colors for labeled items (blue) and citations (red)
\hypersetup{colorlinks=true, linkcolor=blue, citecolor=red}

% Make reference section title font smaller
\renewcommand{\refname}{\large\bf{References}}

% --- Settings for printing computer code ---

% Define colors for green text (comments), grey text (line numbers),
% and green frame around code
\definecolor{greenText}{rgb}{0.5, 0.7, 0.5}
\definecolor{greyText}{rgb}{0.5, 0.5, 0.5}
\definecolor{codeFrame}{rgb}{0.5, 0.7, 0.5}

% Define code settings
\lstdefinestyle{code} {
  frame=single, rulecolor=\color{codeFrame},            % Include a green frame around the code
  numbers=left,                                         % Include line numbers
  numbersep=8pt,                                        % Add space between line numbers and frame
  numberstyle=\tiny\color{greyText},                    % Line number font size (tiny) and color (grey)
  commentstyle=\color{greenText},                       % Put comments in green text
  basicstyle=\linespread{1.1}\ttfamily\footnotesize,    % Set code line spacing
  keywordstyle=\ttfamily\footnotesize,                  % No special formatting for keywords
  showstringspaces=false,                               % No marks for spaces
  xleftmargin=1.95em,                                   % Align code frame with main text
  framexleftmargin=1.6em,                               % Extend frame left margin to include line numbers
  breaklines=true,                                      % Wrap long lines of code
  postbreak=\mbox{\textcolor{greenText}{$\hookrightarrow$}\space} % Mark wrapped lines with an arrow
}

% Set all code listings to be styled with the above settings
\lstset{style=code}

% --- Math/Statistics commands ---

% Add a reference number to a single line of a multi-line equation
% Usage: "\numberthis\label{labelNameHere}" in an align or gather environment
\newcommand\numberthis{\addtocounter{equation}{1}\tag{\theequation}}

% Shortcut for bold text in math mode, e.g. $\b{X}$
\let\b\mathbf

% Shortcut for bold Greek letters, e.g. $\bg{\beta}$
\let\bg\boldsymbol

% Shortcut for calligraphic script, e.g. %\mc{M}$
\let\mc\mathcal

% \mathscr{(letter here)} is sometimes used to denote vector spaces
\usepackage[mathscr]{euscript}

% Convergence: right arrow with optional text on top
% E.g. $\converge[w]$ for weak convergence
\newcommand{\converge}[1][]{\xrightarrow{#1}}

% Normal distribution: arguments are the mean and variance
% E.g. $\normal{\mu}{\sigma}$
\newcommand{\normal}[2]{\mathcal{N}\left(#1,#2\right)}

% Uniform distribution: arguments are the left and right endpoints
% E.g. $\unif{0}{1}$
\newcommand{\unif}[2]{\text{Uniform}(#1,#2)}

% Independent and identically distributed random variables
% E.g. $ X_1,...,X_n \iid \normal{0}{1}$
\newcommand{\iid}{\stackrel{\smash{\text{iid}}}{\sim}}

% Equality: equals sign with optional text on top
% E.g. $X \equals[d] Y$ for equality in distribution
\newcommand{\equals}[1][]{\stackrel{\smash{#1}}{=}}

% Math mode symbols for common sets and spaces. Example usage: $\R$
\newcommand{\R}{\mathbb{R}}   % Real numbers
\newcommand{\C}{\mathbb{C}}   % Complex numbers
\newcommand{\Q}{\mathbb{Q}}   % Rational numbers
\newcommand{\Z}{\mathbb{Z}}   % Integers
\newcommand{\N}{\mathbb{N}}   % Natural numbers
\newcommand{\F}{\mathcal{F}}  % Calligraphic F for a sigma algebra
\newcommand{\El}{\mathcal{L}} % Calligraphic L, e.g. for L^p spaces

% Math mode symbols for probability
\newcommand{\pr}{\mathbb{P}}    % Probability measure
\newcommand{\E}{\mathbb{E}}     % Expectation, e.g. $\E(X)$
\newcommand{\var}{\text{Var}}   % Variance, e.g. $\var(X)$
\newcommand{\cov}{\text{Cov}}   % Covariance, e.g. $\cov(X,Y)$
\newcommand{\corr}{\text{Corr}} % Correlation, e.g. $\corr(X,Y)$
\newcommand{\B}{\mathcal{B}}    % Borel sigma-algebra

% Other miscellaneous symbols
\newcommand{\tth}{\text{th}}	% Non-italicized 'th', e.g. $n^\tth$
\newcommand{\Oh}{\mathcal{O}}	% Big-O notation, e.g. $\O(n)$
\newcommand{\1}{\mathds{1}}	% Indicator function, e.g. $\1_A$

% Additional commands for math mode
\DeclareMathOperator*{\argmax}{argmax}    % Argmax, e.g. $\argmax_{x\in[0,1]} f(x)$
\DeclareMathOperator*{\argmin}{argmin}    % Argmin, e.g. $\argmin_{x\in[0,1]} f(x)$
\DeclareMathOperator*{\spann}{Span}       % Span, e.g. $\spann\{X_1,...,X_n\}$
\DeclareMathOperator*{\bias}{Bias}        % Bias, e.g. $\bias(\hat\theta)$
\DeclareMathOperator*{\ran}{ran}          % Range of an operator, e.g. $\ran(T) 
\DeclareMathOperator*{\dv}{d\!}           % Non-italicized 'with respect to', e.g. $\int f(x) \dv x$
\DeclareMathOperator*{\diag}{diag}        % Diagonal of a matrix, e.g. $\diag(M)$
\DeclareMathOperator*{\trace}{trace}      % Trace of a matrix, e.g. $\trace(M)$

% Numbered theorem, lemma, etc. settings - e.g., a definition, lemma, and theorem appearing in that 
% order in Section 2 will be numbered Definition 2.1, Lemma 2.2, Theorem 2.3. 
% Example usage: \begin{theorem}[Name of theorem] Theorem statement \end{theorem}
\theoremstyle{definition}
\newtheorem{theorem}{Theorem}[section]
\newtheorem{proposition}[theorem]{Proposition}
\newtheorem{lemma}[theorem]{Lemma}
\newtheorem{corollary}[theorem]{Corollary}
\newtheorem{definition}[theorem]{Definition}
\newtheorem{example}[theorem]{Example}
\newtheorem{remark}[theorem]{Remark}

% Un-numbered theorem, lemma, etc. settings
% Example usage: \begin{lemma*}[Name of lemma] Lemma statement \end{lemma*}
\newtheorem*{theorem*}{Theorem}
\newtheorem*{proposition*}{Proposition}
\newtheorem*{lemma*}{Lemma}
\newtheorem*{corollary*}{Corollary}
\newtheorem*{definition*}{Definition}
\newtheorem*{example*}{Example}
\newtheorem*{remark*}{Remark}
\newtheorem*{claim}{Claim}
\newcommand{\Tau}{\mathcal{T}}
\newcommand{\PSet}{\mathcal{P}}


% --- Left/right header text (to appear on every page) ---

% Include a line underneath the header, no footer line
\pagestyle{fancy}
\renewcommand{\footrulewidth}{0pt}
\renewcommand{\headrulewidth}{0.4pt}

% Left header text: course name/assignment number
\lhead{MAM3000W, 3MS -- Hand-in 4}

% Right header text: your name
\rhead{Alex Cristaudo, CRSALE010}

% --- Document starts here ---

\begin{document}
\subsection*{Question 1}

\begin{enumerate}[label=(\alph*)]
\item False. Consider metric space $(\R, d_E)$ and define $A_i = [i, i+1]$. Since each $A_i$ is closed and bounded, it is compact. But then $\bigcup _{n \in \Z} A_n = \R$ which is not bounded under $(\R, d_E)$ and thus not compact

\item True.  Take some $n \in I$. Then since $A_n$ is compact, $A_n$ is totally bounded. Let $\epsilon > 0$. Then there exists some finite set $\{x_1, x_2, ..., x_n\} \subseteq X$ such that $A_n \subseteq \bigcup _{i=1} ^n B(x_i; \epsilon)$. But $\bigcap _{i \in I}A_i \subseteq A_n \subseteq \bigcup _{i=1} ^n B(x_i; \epsilon)$. So $\bigcap _{i \in I}A_i$ is totally bounded. Now since each $A_i$ is complete (follows from $A_i$ being compact),  $\bigcap _{i \in I}A_i$ is also complete. This follows from if $(x_n)$ is Cauchy in $\bigcap _{i \in I}A_i$, then $(x_n)$ is Cauchy in each $A_i$ and thus converges to some $x \in A_i$ for all $i \in I$, so $x \in \bigcap _{i \in I}A_i$. So  $\bigcap _{i \in I}A_i$ is totally bounded and complete, thus is compact

\item False. Consider $f: ([1, \infty), d_E) \to (\R, d_E)$ defined by $f(x) = \frac{1}{x}$. First note that if $x_n \to x$, for $x_n, x \in [1, \infty)$. Then $|f(x_n) - f(x)| = |\frac{1}{x_n} - \frac{1}{x}| = |\frac{x_n - x}{xx_n}| \le |x_n - x| \to 0$ as $n \to \infty$. So $f$ is continuous. But then $[0, 1]$ is compact under $(\R, d_E)$, and $f^{-1}([0,1]) = [1, \infty)$ which is not compact as it is not bounded

\item False. Consider metric space $(\R, d_E)$ with $A_i = \{i\}$. Each $A_i$ is connected. If $I = [-1,1] \cup [2,3]$, then $\bigcup _{i \in I}A_i = I$ which is disconnected so it is not connected

\item False. Consider metric space $(\R ^2, d_E)$. Let $A_1 = \{(x,y) \in \R ^2 : x^2 + y^2 = 4\}$ and $A_2 = \{(x,y) \in \R ^2 : (x-1)^2 + (y-1)^2 = 1\}$. Now consider function $f: [0, 2\pi) \to \R^2$ defined by $f(\theta) = (2\cos \theta , 2\sin \theta)$. This is continuous because each component is continuous. This is a function on interval $[0, 2 \pi )$ which is connected. Then $f([0, 2 \pi)) = A_1$ so $A_1$ is connected. Similarly, we can define $g: [0, 2\pi) \to \R^2$ by $g(\theta) = (1+\cos \theta  , 1+\sin \theta )$ which is also continuous to get that $A_2 = g([0, 2 \pi ))$ is connected. But $A_1 \cap A_2 $ is a set of 2 distinct points, say $A_1 \cap A_2 = \{(x_1, y_1), (x_2, y_2)\}$. This is then disconnected, as $\{(x_1, y_1)\}$ and $\{(x_2, y_2)\}$ are disjoint and closed and their union is $A_1 \cap A_2$, and thus not connected

\end{enumerate}
\newpage
\subsection*{Question 2}
$(\Rightarrow)$ Suppose $(X,d)$ is compact. Let $(A_n)$ be a decreasing sequence of non-empty subsets of $X$. Now let $x_n$ be some value in $A_n$ for each $n \in \N$. Since $(X,d)$ is compact, there exists some subsequence $(x_{n_k})$ such that $x_{n_k} \to x \in X$. Notice that since $x_n \in A_1$ for all $n \in \N$, we get that $x_{n_k}$ is a sequence in $A_1$. Now let $N \in \N$. Now notice that sequence $(x_{n_k + N})$ is a subsequence of $(x_{n_k})$, and $n_k + N \ge k + N \ge N$, and since $x_N \in A_N$, and that $x_n \in A_N$ for all $n \ge N$, $(x_{n_k + N})$ is a sequence in $A_N$. But $x_{n_k + N} \to x$ as well (since it is a subsequence of $(x_{n_k})$), so $x \in \overline{A_N}$. Since this is true for every $N \in \N$, we get $x \in \bigcap _{n \in \N}\overline{A_n}$. \\ 
$(\Leftarrow)$ Suppose that for any decreasing sequence $(A_n)$ of non-empty subsets of $X$, then $\bigcap _{n \in \N} \overline{A_n} \ne \emptyset$. Now let $(x_n)$ be a sequence in $X$. Let $A_n = \{x_k : k \ge n\}$. This is a decreasing sequence of non-empty subsets of $X$. So there exists some $x \in \bigcap _{n \in \N} \overline{A_n}$. Then $x \in \overline{A_N}$ for each $N \in \N$. That means there exists some sequence $(y_{Nn})_{n=1}^\infty$ in $A_N$ such that $y_{Nn} \to x$. We create subsequence $(a_N)$ as follows: For each $N \in \N$, let $a_N = y_{Nn}$ for some $y_{Nn}$ such that $d(y_{Nn}, x) < \frac{1}{N}$. Then $(a_N)$ is a subsequence of $(x_n)$, and $d(a_N, x) < \frac{1}{N}$ for all $N \in \N$, and since $\frac{1}{N} \to 0$ as $N \to \infty$, we get that $(a_n)$ is a convergent subsequence of $(x_n)$, converging to some $x \in X$. Thus $(X,d)$ is compact

\end{document}